\chapter{نتايج}
\section{مقدمه}
ارائه‌ي داده‌ها، نتايج و تحليل و تفسير آنها در فصل سوم ارائه مي‌شود. تفسير و تحليل نتايج نبايد بر اساس حدس و گمان باشد، بلكه بايد برمبناي نتايج استخراج شده از پروژه و يا استناد به تحقيقات ديگران باشد. در ارائه‌ي نتايج با توجه به راهنماي كلي نگارش فصل ها، تا حد امكان تركيبي از نمودار و جدول استفاده شود. فصل نتایج نیز حداکثر شامل 20 صفحه است.

\section{محتوا}
\section{پيشنهادها}
عناوين و موضوعات پيشنهادي را براي تحقيقات آتي بيشتر در زمينه‌ي مورد بحث در آينده ارائه مي‌کند.

\vspace{1cm}
\subsection*{راهنمای ساختاری نگارش}

\begin{table}[h!]
\centering
\caption{بالانويس جدول}
\label{tab:sample_ch3}
\begin{tabular}{cc}
\toprule
\textbf{جدول} & \textbf{متن} \\
\midrule
نمونه داده & توضیحات \\
\bottomrule
\end{tabular}
\end{table}

\begin{figure}[h!]
\centering
\includegraphics[width=0.4\textwidth]{example-image}
\caption{زيرنويس شکل}
\label{fig:sample_ch3}
\end{figure}

\noindent
\textbf{شکل}\\
متن

\vspace{0.5cm}
\noindent
\textbf{فرمول:}
\begin{equation}
f(x) = \int_{0}^{X} \mu \, dD
\end{equation}
متن